\documentclass[a4paper]{article}

\usepackage[english]{babel}
%\usepackage[utf8x]{inputenc}
\usepackage{amsmath}
\usepackage{graphicx}
\usepackage[colorinlistoftodos]{todonotes}
\usepackage[all]{xy}
\usepackage{dsfont}
\usepackage{amssymb}
\usepackage{bbold}
\usepackage{amsthm}

\usepackage{tikz}
\usepackage{verbatim}
\usetikzlibrary{matrix}
\usetikzlibrary{cd}

\usepackage[autostyle=true]{csquotes}
\usepackage{pdfpages}
\usepackage{hyperref}
\usepackage{fancyhdr}



\newtheorem{theorem}{Theorem}[section]
\newtheorem{proposition}{Proposition}[theorem]
\newtheorem{lemma}{Lemma}[theorem]
\newtheorem{corollary}{Corollary}[theorem]
\newtheorem*{definition}{Definition}


\newcommand{\norm}[1]{\left\lVert#1\right\rVert}

\newcommand{\Rn}{\mathbb{R}^n}


\newcommand{\too}[1]{\overset{#1}{\longrightarrow}}
\newcommand{\hookar}[1]{\overset{#1}{\hookrightarrow}}
\newcommand{\colim}[1]{ \underset{#1}{\texttt{colim}} \text{ } }
\newcommand{\fset}[1]{ \{ #1_0 , \dots , #1_n \} }
\newcommand{\ddt}{\frac{d}{dt}}


\newcommand{\N}{\mathbb{N}}
\newcommand{\Z}{\mathbb{Z}}







\fancyhead[L]{\includegraphics[scale=0.3]{UFT.png}}

\begin{document}


\thispagestyle{fancy}
\vspace*{1cm}

\begin{minipage}[r]{12cm}
    \begin{flushright}
\begin{tikzpicture}
\draw[green, very thick] (0,0) rectangle (3,2);
\end{tikzpicture}
    \end{flushright}
\end{minipage}

\begin{center}
\Large{\textbf{Control 2\\ Álgebra Lineal\\ }}
\end{center}
\vspace{1cm}
\begin{minipage}[r]{15cm}
\begin{flushleft}

Fecha: 29 de mayo de 2023 \\%
Duración: 120 minutos  \\
Profesores: Sebastián Cea.\\\hspace{1.7cm} Camila Muñoz\\
\end{flushleft}
\end{minipage}
\vspace{2cm}
\begin{center}
\Large{\textbf{Facultad de Ingeniería}}
\end{center}
\vspace{1cm}
\textbf{En caso de que un alumno sea sorprendido en actitudes deshonestas, ya sea intentando o
dejándose copiar, será calificado con nota mínima 1.0 y perderá toda posibilidad de eliminar
dicha nota. Además, esta conducta quedará archivada como antecedente en su carpeta y
podrá ser causal de eliminación.
Así mismo, está prohibido salir de la sala, usar o mantener sobre la mesa teléfonos celulares
o cualquier tipo de aparato electrónico, a no ser que el profesor lo requiera.}\\



\vspace{5cm}
\makebox[12cm]{\textbf{Nombre Alumno}:\ \hrulefill} 
\newpage
\thispagestyle{empty}
\section*{Preguntas}
\begin{enumerate}


    \item Sean $v = (1, -2, k)$ , $u = (3, k, -2)$ y $ w = (2, -1, 1)$ 
    \begin{enumerate}
        \item {\bf (6 puntos)} ¿Para que valores de $k\in \mathbb{R}$ el vector $v$ es combinaci\'on lineal de los vectores $u$ y $w$?
        \item {\bf (6 puntos)} Si $k=2$, ¿cu\'al es la dimensi\'on del espacio generado?
    \end{enumerate}

    
\newpage

\newpage
\item {\bf (12 puntos)} Determine si el siguiente conjunto es linealmente dependiente o independiente en el $\mathbb{R}-$ espacio vectorial $M_{2}(\mathbb{R})

\[\left\{\begin{pmatrix} -2 & 1 \\ 1 & 0 \end{pmatrix},\begin{pmatrix} 1 & -1 \\ 0 & 2 \end{pmatrix},\begin{pmatrix} -1 & 0 \\ 1 & 2 \end{pmatrix},\begin{pmatrix} 0 & -3 \\ 1 & 2 \end{pmatrix}\right\}. \]

\newpage
    \item Consideremos $V$ un espacio vectorial de dimensión $3$.
    \begin{enumerate}
        \item {\bf (6 puntos)} Sea $\{ v_1, v_2 , v_3 \}$ una base de $V$. Demuestre que \textbf{para todo} $\lambda \neq 0$ el conjunto $\{ \lambda v_1, v_2 , v_3  \}$ es tambi\'en una base.
        
        \item {\bf (6 puntos)} Demuestre que si $\lambda =0$, el conjunto $\{ \lambda v_1, v_2 , v_3  \}$ no es una base de $V$. (\textbf{Hint:} Intente escribir el vector $v_1$ como combinaci\'on lineal de los vectores $\lambda v_1, v_2 , v_3 $).

    \end{enumerate}

    \newpage
\newpage

    \item Sea $H$ el subconjunto de $\mathbb{R}^3$ definido por
\[H=\{(x,y,z)\in\mathbb{R}^3\colon x+y-2z=0  \}\]

\begin{enumerate}
    \item {\bf (4 puntos)} Demuestre que $H$ es un subespacio vectorial de $\mathbb{R}^3$.
    \item {\bf (4 puntos)} Encuentre una base de $H$ y determine su dimensión. 
    \item {\bf (4 puntos)} ?`Es una base ortogonal??`Es ortonormal? En caso de no serlo, ortonormalice la base.
\end{enumerate}

        
        
   \newpage 

\section*{Formulario}

\begin{itemize}
    \item {\bf Subespacio vectorial: }Si para todo $\alpha \in K$ y para todo $x, y \in H$, $\alpha x+y \in H$, entonces $H$ es subespacio vectorial.
    \item {\bf Combinación lineal: } Un vector de la forma $v = a_1 v_1 + a_2v_2 + \cdots +a_n v_n$, donde $a_1, a_2, \dots a_n$ son escalares, se denomina una combinación lineal de $v_1, v_2, \dots , v_n$.
\item {\bf Linealmente Dependiente e Independiente: }El conjunto $A=\{v_1, v_2, \dots, v_n\}$ es
Linealmente Independiente (L.I.) si $ a_1 v_1 + a_2v_2 + \cdots +a_n v_n=0$ implica que $a_i=0$ para todo $i=1, \dots, n.$ En caso contrario, es linealmente dependiente (L.D.).
\item {\bf Base: } Un conjunto $\{ v_1,v_2, \dots, v_n \}$ es una base para un espacio vectorial $V$ si $\{ v_1,v_2, \dots, v_n \}$ es linealmente independiente y genera a $V$.
\item \textbf{Procedimiento para obtener una base ortonormal}
Sea $B = \{v_i\}_{i=1}^m $ una base ordenada cualquiera de $W $. 
\begin{itemize}
    \item Escribimos $w_1 = v_1$. 
    \item Calculamos $w_2=v_2-\dfrac{<v_2,w_1>}{<w_1,w_1>}w_1$.
    \item Calculamos $w_3=v_3-\dfrac{<v_3,w_1>}{<w_1,w_1>}v_1- \dfrac{<v_3,w_2>}{<w_2,w_2>}w_2$.
    \item Inductivamente obtenemos los vectores $w_m=v_m-\sum_{i=1}^{m-1}c_iw_i$ con $c_i=\frac{<v_m,w_i>}{<w_i,w_i>}$.
\end{itemize}
El conjunto $\{ w_1,w_2, \dots, w_n \}$ es una base ortogonal. Si se normaliza cada vector $w_k$, obtenemos nuevos
vectores $u_k=\frac{w_k}{\|w_k\|}$. Luego, $\{ u_1,u_2, \dots, u_n \}$ una base ortonormal para $W$.
\end{itemize}

\end{itemize}
\end{enumerate}



\end{enumerate}


\end{document}